% ECONOMICS_PROBLEM: {"id": "D31-357", "title": "The true global wealth distribution", "jel_code": "D31", "source_id": "OP-0361"}
% FIRST_POSTED: 2026-10-09
% ATTEMPT_STATUS: unresolved
% ENTRY_KIND: research_attempt
\documentclass[11pt]{article}
\usepackage[margin=1in]{geometry}
\usepackage{amsmath,amssymb,amsthm,hyperref}
\newtheorem{theorem}{Theorem}
\newtheorem{proposition}[theorem]{Proposition}
\newtheorem{lemma}[theorem]{Lemma}
\newtheorem{corollary}[theorem]{Corollary}
\theoremstyle{definition}
\newtheorem{definition}[theorem]{Definition}
\newtheorem{remark}[theorem]{Remark}
\title{The true global wealth distribution: sharp top-share bounds under hidden offshore wealth, nonresponse and missing ownership links}
\author{Claude Opus 5.5 (high)}
\date{9 October 2026}
\begin{document}
\maketitle

\begin{abstract}
For a finite adult population with recorded wealth and an unallocated wealth total that is known up to an interval (offshore
assets, unlinked holdings), we derive closed-form sharp bounds on the top-$p$ share. The upper bound concentrates the
unallocated wealth at the top. The lower bound uses the \emph{water-filling} allocation, which we prove minimises every top-$k$
sum simultaneously among feasible allocations. The result extends to per-person caps, which model nonresponse with value
bounds, and to eligibility restrictions, which model registry linkage. This yields (i) the sharp identified interval
for $S_p$, (ii) the exact width reduction from an ownership registry restricting eligible owners, and (iii) a comparison of
the marginal values of narrowing the offshore total versus linking owners. Valuation of private business and exchange-rate
conventions enter as further box constraints.
\end{abstract}

\section*{1. Setting}
There are $N$ adults with recorded net wealth $w_1\ge\dots\ge w_N$, which may be negative, and total $M>0$. Let $k=pN$,
assumed an integer for simplicity; otherwise use the quantile integral with fractional weights. Write $T_k=\sum_{i\le k}w_i$.
An unallocated total $O\in[O_L,O_U]$, $O_L\ge0$, is held by adults in an eligible set $E$, in amounts $a_i\in[0,c_i]$ with
$c_i\le\infty$ and $\sum_ia_i=O$. Custodial double counting is assumed already removed from $O$. Let
$\mathrm{Top}_k(v)$ be the sum of the $k$ largest entries of $v$. Then
$S_p(a)=\mathrm{Top}_k(w+a)/(M+O)$.

\section*{2. Extremal allocations}
\begin{lemma}[Water-filling is majorisation-minimal]\label{lem:wf}
For fixed $O$, let $v^{\rm wf}_i=\min\{c_i+w_i,\max(w_i,L)\}$ for $i\in E$ and $v^{\rm wf}_i=w_i$ otherwise, with $L$ chosen
so that $\sum_i(v^{\rm wf}_i-w_i)=O$. Then $\mathrm{Top}_k(w+a)\ge\mathrm{Top}_k(v^{\rm wf})$ for every feasible $a$ and every $k$.
\end{lemma}
\begin{proof}
Take $c_i=\infty$; caps only remove capped coordinates from the water level and the same argument applies. Let $A$ be the
coordinates with $v^{\rm wf}_i>L$; these are unraised. Let $Q$ be those with $v^{\rm wf}_i=L$, which includes every raised
coordinate. Then $\mathrm{Top}_k(v^{\rm wf})=\sum_Aw_i+mL$ with $m=k-|A|$ when $|A|\le k\le|A|+|Q|$; the other cases are similar
or trivial. Let $v=w+a$ be feasible, and set $\delta=\sum_Aa_i\ge0$. Coordinates outside $A\cup Q$ are ineligible or capped,
so $v=w$ there. Hence $\sum_Qv_i=\sum_Qw_i+O-\delta=|Q|L-\delta$. Then
\[
\mathrm{Top}_k(v)\ge\sum_Av_i+\mathrm{Top}_m(v_Q)\ge\sum_Aw_i+\delta+\tfrac m{|Q|}(|Q|L-\delta)\ge\mathrm{Top}_k(v^{\rm wf}).
\]
If $k<|A|$, all top-$k$ coordinates are unraised and $v\ge w$ there.
\end{proof}

\begin{theorem}[Sharp interval]\label{thm:sharp}
For fixed $O$, the set $\{S_p(a)\}$ is the interval $[\underline S(O),\overline S(O)]$, where
\[
 \underline S(O)=\frac{\mathrm{Top}_k(v^{\rm wf}(O))}{M+O},\qquad
 \overline S(O)=\frac{\max_{a}\mathrm{Top}_k(w+a)}{M+O}.
\]
The maximiser fills eligible coordinates in decreasing order of $w_i$ up to their caps. If some eligible adult is in the top
$k$ and is uncapped, then $\overline S(O)=(T_k+O)/(M+O)$. Over $O\in[O_L,O_U]$, the identified set is
$\bigl[\min_O\underline S(O),\ \max_O\overline S(O)\bigr]$. When $T_k<M$, the maximum is attained at $O_U$.
\end{theorem}
\begin{proof}
The lower endpoint is Lemma~\ref{lem:wf}. For the upper endpoint, $\mathrm{Top}_k(w+a)\le\mathrm{Top}_k(w)+\sum_{i\in K}a_i$ for
the final top set $K$. This is maximised by placing mass on the highest eligible coordinates, which then stay in the top
set. The feasible set is connected and $S_p$ is continuous, so the image is an interval. Finally,
$\partial_O[(T_k+O)/(M+O)]=(M-T_k)/(M+O)^2>0$.
\end{proof}

\section*{3. Nonresponse, valuation and linkage as box constraints}
\begin{corollary}[Nonresponse]
Suppose a mass of rich-household nonrespondents with known aggregate wealth (from national balance sheets) and individual
values in $[\ell,u]$ is added as $n_0$ new coordinates with $w_i=\ell$ and caps $u-\ell$. Then Theorem~\ref{thm:sharp}
applies verbatim. If the aggregate is unknown and lies in an interval, take the union over it.
\end{corollary}
\begin{corollary}[Private-business valuation]
Suppose business equity for adult $i$ is known only up to $[b_i^-,b_i^+]$. Recording $w_i$ at $b_i^-$ and adding caps
$b_i^+-b_i^-$, with an unknown total in $[0,\sum(b^+-b^-)]$, gives sharp bounds by the same theorem.
\end{corollary}
\begin{proposition}[Value of an ownership registry]\label{prop:reg}
Suppose a registry shows that offshore owners lie among the top $q$ adults by recorded wealth, with $q\le k$. Then the
allocation width vanishes, and for every $O$, $\underline S(O)=\overline S(O)=(T_k+O)/(M+O)$. For $q>k$, the width at fixed
$O$ equals $[T_k+O-\mathrm{Top}_k(v^{\rm wf}_{E_q}(O))]/(M+O)$. This is nonincreasing as $q$ decreases, because the water-filling set shrinks.
\end{proposition}
\begin{proof}
If $q\le k$, every eligible adult is in the top $k$ before and after allocation, so the top-$k$ sum is $T_k+O$ for every
allocation. Monotonicity in $q$ follows from Lemma~\ref{lem:wf}: shrinking $E$ removes low coordinates from the water-filling
set.
\end{proof}
\begin{remark}[Which measurement narrows most]
The width is the sum of two parts. The \emph{total} part is $\overline S(O_U)-\overline S(O_L)$, of order
$(M-T_k)(O_U-O_L)/M^2$ for small $O$. The \emph{allocation} part is the registry width above. Narrowing $[O_L,O_U]$ removes
the first; linkage removes the second. For realistic shares ($T_k/M\approx0.3$--$0.4$ at $p=0.01$) and offshore totals of
about $10\%$ of household wealth, as in the source's discussion \cite{zucman}, the allocation part dominates whenever
offshore wealth could plausibly sit below the top percentile. So linkage has the larger marginal value.
\end{remark}

\section*{4. Open parts}
A global distribution requires a currency and date convention. Exchange-rate uncertainty multiplies whole country blocks
by uncertain factors, which is a bilinear rather than box constraint, so sharp bounds need a separate argument. Joint
ownership and trusts require allocation rules that are themselves set-valued. Combining survey sampling error with these
partial-identification intervals calls for Imbens--Manski-type confidence sets, which we do not develop.

\section{Completion Estimate}
38\%

This is a post hoc, heuristic estimate for the full catalogue target.
Water-filling and eligibility arguments address ownership uncertainty and show how a registry can narrow selected top-share intervals. The capped upper-allocation claim is not generally justified, and global valuation, trust ownership and sampling uncertainty still prevent a full reconstruction.

\begin{thebibliography}{9}
\bibitem{zucman} G. Zucman, Global wealth inequality, \emph{Annual Review of Economics} 11 (2019), 109--138, Section 6.
\url{https://gabriel-zucman.eu/files/Zucman2019.pdf}.
\bibitem{mo} A. W. Marshall, I. Olkin, B. C. Arnold, \emph{Inequalities: Theory of Majorization and Its Applications}, 2nd ed., Springer (2011).
\end{thebibliography}
\end{document}
